\documentclass[11pt, a4paper]{article}
\usepackage[utf8]{inputenc}
\usepackage[T1]{fontenc}
\usepackage[spanish]{babel}
\usepackage{amsmath, amssymb}
\usepackage{geometry}
\usepackage{enumitem}
\usepackage{titlesec}
\usepackage{parskip}

\geometry{margin=2.5cm}

\titleformat{\section}{\large\bfseries}{}{0em}{}
\titleformat{\subsection}{\normalsize\bfseries}{}{0em}{}
\titleformat{\subsubsection}{\normalsize\itshape}{}{0em}{}

\title{\textbf{Resumen -- Neurociencia y Comportamiento} \\
\large Inteligencia Artificial y Neurociencias \\ Universidad Torcuato Di Tella}
\author{}
\date{}

\begin{document}
\maketitle

%% ============================================================
\section{Clase 1 -- Modelado Computacional del Comportamiento Humano}
%% ============================================================

\subsection{Contexto: el paper de van Opheusden et al.\ (2023)}

El art\'iculo \textit{Expertise increases planning depth in human gameplay} (Nature, 2023) estudia la \textbf{profundidad de planificaci\'on} en humanos: \textquestiondown los jugadores m\'as expertos planifican m\'as pasos a futuro que los novatos?

\textbf{Problema de medici\'on.} Medir profundidad de planificaci\'on es dif\'icil porque:
\begin{itemize}[nosep]
    \item La complejidad de muchas tareas no basta para distinguir expertos de novatos (e.g.\ tareas de dos pasos).
    \item El juego debe ser \emph{tratable computacionalmente} (no como el ajedrez).
    \item Debe tener reglas simples, novedosas e interesantes.
\end{itemize}
El juego \textbf{Four-In-A-Row} (4 en l\'inea) cumple todas estas condiciones: es suficientemente complejo para que los expertos planifiquen con profundidad, pero tambi\'en simple para ajustar un modelo computacional con par\'ametros bien definidos.

\subsection{El modelo cognitivo}

El modelo tiene tres componentes, cada uno con par\'ametros ajustables:

\begin{enumerate}[nosep]
    \item \textbf{Feature-based Value Function} (funci\'on heur\'istica de valor): asigna un valor $V(s, \mathbf{w})$ a cada estado del tablero $s$.
    \[
        V(s, \mathbf{w}) = \sum_{i=1}^{5} w_i \, \phi_i(s, \text{self}) \;-\; \sum_{i=1}^{5} w_i \, \phi_i(s, \text{oponente})
    \]
    donde $\phi_i$ son cinco \emph{features}: centro, dos-en-l\'inea conectado, dos-en-l\'inea no conectado, tres-en-l\'inea y cuatro-en-l\'inea. Features del mismo color comparten pesos. Los pesos del jugador activo se multiplican por una constante de escala $C$.

    \item \textbf{Search Algorithm} (algoritmo de b\'usqueda): en cada iteraci\'on elige una posici\'on del \'arbol para expandir, eval\'ua las posiciones resultantes con $V(s)$ y \emph{poda} movimientos cuyo valor est\'e por debajo del mejor menos un umbral. Despu\'es de cada iteraci\'on, el algoritmo se detiene con probabilidad $\gamma$ (distribuci\'on geom\'etrica sobre el n\'umero total de iteraciones). Los nodos rojos del \'arbol de decisi\'on forman la \textbf{variaci\'on principal} (secuencia de mejores movimientos para ambos jugadores).

    \item \textbf{Attention Mechanism} (mecanismo de atenci\'on): modela la atenci\'on selectiva descartando aleatoriamente \emph{features} (en posiciones y orientaciones espec\'ificas) antes de construir el \'arbol. Adem\'as se a\~nade ruido gaussiano a $V(s)$ en cada nodo y existe una tasa de lapsus $\lambda$ (probabilidad de elegir una acci\'on al azar).
\end{enumerate}

\subsection{Resultados principales}

\begin{itemize}[nosep]
    \item El modelo predice con buena precisi\'on las \textbf{elecciones} humanas (validaci\'on cruzada 5-fold en partidas humano vs.\ humano), muy por encima del azar.
    \item Predice \textbf{tiempos de respuesta}: el tama\~no del \'arbol de decisi\'on es predictor del RT observado.
    \item Predice \textbf{movimientos oculares}: la distribuci\'on de casillas visitadas por el algoritmo de b\'usqueda correlaciona con la distribuci\'on de atenci\'on visual ($\rho = 0.535 \pm 0.024$, $P < 0.001$).
\end{itemize}

\textbf{Conclusi\'on:} los jugadores m\'as fuertes planifican m\'as profundamente y tienen menos lapsus atencionales. No hay evidencia de mejora en la calidad de la heur\'istica (pesos de features).

\subsection{Hip\'otesis de implementaci\'on neural}

\begin{center}
\begin{tabular}{ll}
\hline
\textbf{Componente del modelo} & \textbf{Regi\'on cerebral} \\
\hline
Feature-based Value Function & Corteza orbitofrontal \\
Search Algorithm & Corteza prefrontal, c\'ingulo, hipocampo \\
Attention Mechanism & Red frontoparietal \\
\hline
\end{tabular}
\end{center}

\subsection{Discusi\'on: profundidad vs.\ reconocimiento de patrones}

La literatura explica la superioridad de los expertos por mejor \emph{pattern recognition} y/o b\'usqueda m\'as profunda. Los autores equiparan ``planificar m\'as profundamente'' con mayor ``velocidad de procesamiento'': los expertos afinan su representaci\'on de features relevantes, lo que permite m\'as evaluaciones por unidad de tiempo. Esto es consistente con mejor reconocimiento de patrones, pero destaca el rol subestimado de la velocidad de procesamiento.

%% ============================================================
\section{Clase 2 -- Aprendizaje (Parte 1)}
%% ============================================================

\subsection{Definici\'on de aprendizaje}

\textbf{Aprendizaje:} cambio en el comportamiento debido a la experiencia.

Tipos de aprendizaje simple en humanos (condicionamiento cl\'asico de Pavlov y condicionamiento operante de Skinner se ver\'an en otra clase).

\subsection{Aprendizaje en redes neuronales artificiales}

Las neuronas biol\'ogicas (soma, ax\'on, dendritas) inspiran las neuronas artificiales: suma ponderada de entradas + no-linealidad. Las redes neuronales artificiales organizan estas unidades en capas (entrada, ocultas, salida).

\subsection{Aprendizaje en organoides inteligentes}

\begin{itemize}[nosep]
    \item \textit{Brain organoid reservoir computing for AI} (Nature Electronics, 2023): hardware inspirado en la estructura cerebral para reconocimiento de audio.
    \item \textit{In vitro neurons learn and exhibit sentience when embodied in a simulated game-world} (2022): neuronas in vitro aprenden a jugar Pong. La actividad neural cambia en tiempo real para minimizar la impredecibilidad ambiental (\emph{Free Energy Principle}).
\end{itemize}

\subsection{Priors humanos para videojuegos}

Paper: \textit{Investigating Human Priors for Playing Video Games} (Dubey et al., UC Berkeley). Pregunta central: \textquestiondown por qu\'e los humanos aprenden juegos como Mario Bros.\ r\'apidamente y las m\'aquinas tardan much\'isimo m\'as?

\subsection{Aprendizaje humano: \textquestiondown cu\'ales son los priors?}

Ejemplos de conocimiento innato (\emph{priors}) en humanos:

\begin{itemize}[nosep]
    \item \textbf{Detecci\'on de amenazas:} estamos cableados para encontrar v\'iboras m\'as r\'apido que flores.
    \item \textbf{Violaci\'on de leyes f\'isicas:} beb\'es de 3 meses muestran sorpresa ante eventos f\'isicamente imposibles.
    \item \textbf{Intuici\'on num\'erica y probabil\'istica:} los beb\'es miran m\'as si cae una pelota de color menos probable; generalizan entre modalidades (auditiva $\to$ visual).
    \item \textbf{Razonamiento l\'ogico pre-verbal:} beb\'es de 12--19 meses tienen nociones de sumas, restas y conservaci\'on (Cesana-Arlotti et al.).
    \item \textbf{Preferencia por caras:} incluso intrauterinamente, los fetos prefieren est\'imulos tipo cara.
\end{itemize}

Estos comportamientos son \textbf{instintivos}: aparecen en humanos t\'ipicos con desarrollo t\'ipico. Permiten el aprendizaje de comportamientos m\'as complejos (e.g.\ atenci\'on innata al habla $\to$ aprendizaje del lenguaje). Fueron seleccionados a lo largo de la evoluci\'on.

\subsection{Evoluci\'on y selecci\'on natural}

\textbf{Evoluci\'on:} cambio en los organismos de una poblaci\'on a lo largo del tiempo.

\textbf{Selecci\'on natural} (Charles Darwin): ocurre cuando existen tres condiciones:
\begin{enumerate}[nosep]
    \item \textbf{Variabilidad} entre los rasgos de los individuos.
    \item \textbf{Heredabilidad} de esas variaciones (aunque pueden ocurrir mutaciones o recombinaciones).
    \item \textbf{Adaptabilidad:} algunas variaciones se adaptan mejor al ambiente; los individuos que las poseen sobreviven m\'as y dejan m\'as descendencia.
\end{enumerate}

\subsubsection{Herencia: concepto de gen}

Un \textbf{gen} es la unidad m\'inima de herencia: secuencia de ADN que codifica un determinado producto en el organismo. El concepto es anterior al descubrimiento del ADN y surgi\'o de observar que padres e hijos se parecen en aspectos que se heredan independientemente.

\subsubsection{Darwin vs.\ Lamarck}

En la teor\'ia de \textbf{Lamarck}, los rasgos adquiridos durante la vida se heredan, por lo que no se necesitan variaciones previas en la poblaci\'on. En la de \textbf{Darwin}, los rasgos adquiridos \emph{no} se heredan; las variaciones preexistentes son esenciales para la selecci\'on natural.

\subsubsection{Tres frutos de un proceso evolutivo darwinista}

\begin{enumerate}[nosep]
    \item \textbf{Adaptaciones:} rasgos que aumentaron el \'exito reproductivo (e.g.\ cord\'on umbilical, hemoglobina). Ejemplos comportamentales: gusto por el az\'ucar, impulso sexual, apego paterno-materno, enamoramiento.
    \item \textbf{Subproductos (\emph{byproducts}):} consecuencias de rasgos adaptativos sin ventaja evolutiva propia (e.g.\ ombligo, color rojo de la sangre). Ejemplos comportamentales: adicciones, gusto por el f\'utbol, sexo con anticonceptivos.
    \item \textbf{Ruido:} variaci\'on aleatoria sin funci\'on.
\end{enumerate}

\subsection{Selecci\'on de comportamientos y estructura cerebral}

El cerebro es una gran red neuronal cuyos \textbf{hiperpar\'ametros} (estructura a gran escala) fueron seleccionados por evoluci\'on darwinista. Se seleccionan genes que influyen en la estructura cerebral: si una mutaci\'on mejora supervivencia y reproducci\'on, esos genes mutados se seleccionan.

Analog\'ia v\'alida: la selecci\'on natural ``dise\~n\'o'' la estructura a gran escala y los hiperpar\'ametros; durante la vida hacemos el \textbf{\emph{fine tuning}}.

Cita de Stanislas Dehaene: el cerebro no es tabula rasa ni m\'odulos fijos e inmutables. Los circuitos est\'an organizados desde el nacimiento pero poseen plasticidad. Lo innato y lo adquirido se combinan.

%% ============================================================
\section{Clase 3 -- Aprendizaje (Parte 2)}
%% ============================================================

\subsection{Teor\'ia de la Carga Cognitiva (CLT)}

La clave del aprendizaje es guardar informaciones y automatizar procedimientos en la \textbf{memoria de largo plazo}. La mejor forma: ense\~nar y aprender de forma expl\'icita y estructurada.

Ventajas: permite repetici\'on desde el principio (principio de Hebb), evita frustraci\'on, optimiza motivaci\'on y refuerza plasticidad gracias a la dopamina. Resolver problemas solos desde el comienzo (aunque sean sencillos) genera motivaci\'on, tanto en ni\~nos como en adultos.

\subsection{Las tres etapas de cualquier aprendizaje}

\begin{enumerate}[nosep]
    \item \textbf{Adquisici\'on:} el estudiante a\'un no responde con precisi\'on sin ayuda. Instrucci\'on expl\'icita y sistem\'atica. Pr\'actica en bloques (preguntas del mismo tipo).
    \item \textbf{Fluidez:} responde con precisi\'on pero lentamente. Pr\'actica cronometrada para desarrollar velocidad y automaticidad.
    \item \textbf{Generalizaci\'on:} tiene fluidez y transfiere la habilidad a contextos nuevos. Pr\'actica intercalada y problemas no rutinarios.
\end{enumerate}
Adem\'as existe una etapa de \textbf{mantenimiento}: retener la habilidad mediante pr\'actica espaciada.

\subsection{Los cuatro pilares del aprendizaje (Dehaene)}

\subsubsection{1.\ Atenci\'on}

Aprendemos cuando hacemos un esfuerzo que nos saca levemente de la \textbf{zona de confort} $\to$ \textbf{pr\'actica deliberada}.

Concepto de \textbf{OK Plateau}: al aprender a manejar o teclear, se llega a un nivel aceptable y se deja de prestar atenci\'on a c\'omo mejorar. Sin atenci\'on no hay mejora. Para pasar de amateur a experto hay que seguir prestando atenci\'on.

Test de la zona de confort (Kahneman, \textit{Pensar r\'apido, pensar despacio}): si puedo hacer lo que hago y al mismo tiempo otra cosa no automatizada, entonces no estoy prestando atenci\'on.

\subsubsection{2.\ Feedback}

El premio por buen desempe\~no funciona mejor que el castigo por mal desempe\~no.

\textbf{Regresi\'on a la media:} puede generar la ilusi\'on de que premiar empeora y castigar mejora. En realidad, tras un desempe\~no extremo (bueno o malo), lo m\'as probable es volver al promedio.

\textbf{Feedback constructivo} (t\'ecnica organizacional):
\begin{itemize}[nosep]
    \item Valoro...
    \item Me pregunto...
    \item Sugiero...
\end{itemize}

\textbf{La maldici\'on del conocimiento} (Steven Pinker): cuando conocemos algo, dejamos de poder ver el mundo como alguien que no lo conoce. Por eso es fundamental el feedback en la buena comunicaci\'on.

\subsubsection{3.\ Motivaci\'on}

Dos tipos principales:
\begin{itemize}[nosep]
    \item \textbf{Intr\'inseca} (hago por placer): satisfacci\'on inherente, sin recompensa externa.
    \item \textbf{Extr\'inseca} (hago por recompensa): obtener recompensa externa o evitar consecuencia negativa.
\end{itemize}

Amenazas a la motivaci\'on en proyectos de producci\'on por pares (software libre):
\begin{enumerate}[nosep]
    \item Temor a la apropiaci\'on unilateral.
    \item Temor a la falla en la integraci\'on (esfuerzo desperdiciado).
\end{enumerate}

\textbf{T\'ecnica Pomodoro:} regulaci\'on del esfuerzo mediante bloques de trabajo con pausas programadas.

\subsubsection{4.\ Consolidaci\'on}

\textbf{El poder del h\'abito:} el aprendizaje es \textbf{multiplicativo}, no aditivo. Cada nuevo conocimiento se relaciona con todos los anteriores, generando crecimiento geom\'etrico. Esto explica la sensaci\'on de que un experto es un ``genio''.

\textbf{Reserva cognitiva:} tolerancia cognitiva o ps\'iquica frente a cambios cerebrales fisiol\'ogicos relacionados con la edad o alguna patolog\'ia. Las \textbf{Zonas Azules} (Loma Linda, Sardinia, Okinawa, Ikaria, Nicoya) muestran que el aprendizaje continuo y los h\'abitos saludables contribuyen a mayor longevidad y reserva cognitiva.

%% ============================================================
\section{Clase 4 -- Emociones}
%% ============================================================

\subsection{Los dos sistemas de Kahneman}

Daniel Kahneman (\textit{Pensar r\'apido, pensar despacio}) distingue dos modos de pensamiento:
\begin{itemize}[nosep]
    \item \textbf{Sistema 1:} r\'apido, autom\'atico, intuitivo, emocional, sin esfuerzo. Es el sistema \emph{emocional}.
    \item \textbf{Sistema 2:} lento, deliberado, racional, con esfuerzo y control voluntario.
\end{itemize}

\subsection{\textquestiondown Qu\'e es una emoci\'on?}

Las emociones son \textbf{estructuras innatas} (los ``hiperpar\'ametros'' seleccionados por la evoluci\'on) que preparan al organismo para responder de forma adaptativa a ciertas situaciones. Son \textbf{adaptativas}: fueron perfeccionadas como dispositivos de supervivencia.

\textbf{Sentimientos vs.\ emociones:} la \emph{emoci\'on} es la respuesta corporal/fisiol\'ogica autom\'atica; el \emph{sentimiento} es la experiencia consciente y subjetiva de esa emoci\'on.

\subsection{Cuatro afirmaciones sobre las que la psicolog\'ia evolutiva s\'i concuerda}

\begin{enumerate}[nosep]
    \item Las emociones son resultado del \textbf{perfeccionamiento de dispositivos de supervivencia} a lo largo de la evoluci\'on de las especies.
    \item Esos dispositivos usan las \textbf{mismas estructuras cerebrales y los mismos neurotransmisores en todos los mam\'iferos}, incluido el ser humano.
    \item Esas estructuras est\'an vinculadas a \textbf{funciones b\'asicas de regulaci\'on del equilibrio interno} necesario para la supervivencia.
    \item Las emociones est\'an arraigadas en fen\'omenos sensitivos viscerales corporales que llamamos \textbf{interocepci\'on} (percepci\'on de estados viscerales: pulsaciones, respiraci\'on, tensi\'on, est\'omago). En la formulaci\'on extrema de \textbf{William James}, ``los cambios f\'isicos que experimentamos ante una situaci\'on \emph{son} las emociones''.
\end{enumerate}

\subsection{Emociones b\'asicas (o primarias)}

Se las llama b\'asicas por tres razones:
\begin{itemize}[nosep]
    \item Aparecen desde \textbf{edades muy tempranas} en todos los humanos (incluso en personas con ceguera y sordera cong\'enitas).
    \item Son \textbf{universales} (en todas las culturas; idea que ya propon\'ia Darwin).
    \item Se reconocen a partir de \textbf{expresiones faciales universales}, independientes del aprendizaje cultural.
\end{itemize}
Ejemplos: miedo, asco, ira, tristeza, alegr\'ia, sorpresa. Ventaja evolutiva: el \textbf{asco} protege de ingerir sustancias t\'oxicas o en mal estado; el \textbf{miedo} activa el organismo para pelear o huir (lucha o huida) ante una amenaza.

\subsection{El sistema nervioso}

\begin{itemize}[nosep]
    \item \textbf{Sistema nervioso central} (cerebro y m\'edula espinal) y \textbf{perif\'erico} (nervios).
    \item Sistema nervioso perif\'erico: \textbf{som\'atico} (movimiento voluntario) y \textbf{aut\'onomo} (ANS, funciones involuntarias).
    \item El aut\'onomo se divide en \textbf{simp\'atico} (SNS, activaci\'on: lucha o huida) y \textbf{parasimp\'atico} (PNS, calma: descanso y digesti\'on).
\end{itemize}

\subsection{Dolor y sufrimiento: regulaci\'on emocional}

\textbf{Dolor $\to$ Sufrimiento:} el dolor (f\'isico o emocional) se convierte en sufrimiento a trav\'es del \textbf{pensamiento terribilizador} (amplificar mentalmente lo negativo).

\textbf{Segundos dardos} (Rick Hanson, \textit{Buddha's Brain}): el primer dardo es el est\'imulo doloroso inevitable; los ``segundos dardos'' son las reacciones que agregamos nosotros (v\'ia SNS, am\'igdala, eje HPAA), y son evitables.

\textbf{Sapolsky} (\textit{\textquestiondown Por qu\'e las cebras no tienen \'ulcera?}): las cebras activan el estr\'es solo ante la amenaza real y luego se relajan; los humanos mantenemos el estr\'es activado con pensamientos, lo que enferma (\'ulcera, etc.).

\textbf{Regulaci\'on emocional como t\'ecnica de mantener a raya el ANS:} calmar el ANS estimulando el PNS orienta cuerpo, cerebro y mente hacia el bienestar. Buscar y \textbf{absorber experiencias positivas} (convertir hechos positivos en experiencias, saborearlas, integrarlas) construye estructura neuronal (neuroplasticidad). Existe una ``receta parasimp\'atica basal'' para orientar el cuerpo hacia la calma.

\textbf{Visi\'on alternativa (David Foster Wallace, \textit{Esto es agua}):} el manejo emocional como capacidad de elegir conscientemente en qu\'e prestar atenci\'on y c\'omo pensar; esa es la verdadera libertad.

%% ============================================================
\section{Clases 5 y 6 -- Gen\'etica del Comportamiento}
%% ============================================================

\subsection{\textquestiondown Qu\'e es la gen\'etica del comportamiento?}

Ciencia que estudia las \textbf{variaciones (diferencias) entre personas}: en qu\'e medida las diferencias en el comportamiento se deben a diferencias en los \textbf{genes} o a diferencias en los \textbf{ambientes}. Estudia el \textbf{componente gen\'etico} (o heredabilidad) del comportamiento. Rasgos estudiados: fisiol\'ogicos (altura, obesidad, estructura cerebral), capacidades cognitivas (memoria, fluencia verbal, matem\'atica), personalidad (extroversi\'on, neuroticismo) y psiqui\'atricos (esquizofrenia, autismo, depresi\'on).

\subsection{Ojo con la palabra ``gen\'etica'': dos definiciones contradictorias}

\begin{itemize}[nosep]
    \item \textbf{En psicolog\'ia evolutiva:} explicaci\'on sobre lo que \emph{todos tenemos en com\'un} (interacci\'on entre ambientes ``t\'ipicos'' y genes ``t\'ipicos''; e.g.\ habilidad de aprender lenguaje, gusto por el az\'ucar).
    \item \textbf{En gen\'etica del comportamiento:} explicaci\'on sobre lo que tenemos \emph{de diferente}, las \emph{variaciones individuales}. En el comportamiento humano influyen muchos genes: hablamos de una \emph{influencia} (componente gen\'etico), no una determinaci\'on.
\end{itemize}

\subsection{Definici\'on rigurosa de componente gen\'etico}

El \textbf{componente gen\'etico $(C_G)$} de un rasgo es la variaci\'on entre individuos debida a variaciones en los genes, dividida por la variaci\'on total:
\[
    C_G = \frac{V_G}{V_G + V_A} = \frac{V_G}{V_T}
\]
donde $V_G$ es la varianza gen\'etica y $V_A$ la ambiental. \textbf{Importante:} el componente gen\'etico \textbf{vale para una dada poblaci\'on en un dado momento del tiempo}; no es una propiedad fija del rasgo. Por definici\'on, componente gen\'etico + componente ambiental $= 100\%$.

\subsection{Componente ambiental: compartido vs.\ no compartido}

El componente ambiental se divide en dos grupos: $V_A = V_{AC} + V_{ANC}$.
\begin{itemize}[nosep]
    \item \textbf{Ambiente compartido (AC):} todo lo que no es gen\'etico y hace que miembros de una familia se parezcan (clase social, estilo parental, barrio). Lo que comparten dos hermanos criados juntos pero no dos separados al nacer.
    \item \textbf{Ambiente no compartido (ANC):} todo lo que no es gen\'etico ni se comparte entre miembros de una familia; ligado a la experiencia particular de cada individuo.
\end{itemize}

Valores t\'ipicos de $C_G$: autismo $\approx 80\%$, esquizofrenia $\approx 50\%$, depresi\'on $37$--$48\%$, neuroticismo $\approx 36\%$, extroversi\'on $\approx 37\%$, zurdera $\approx 25\%$, IMC $40$--$80\%$. Estructuras cerebrales: volumen total del cerebro $\approx 82\%$.

\subsection{M\'etodos cl\'asicos de medici\'on}

\subsubsection{Comparaci\'on gemelos vs.\ mellizos}

Los \textbf{gemelos} (monocig\'oticos) comparten $\approx 100\%$ de sus alelos; los \textbf{mellizos} (dicig\'oticos), $\approx 50\%$. Si la correlaci\'on de un rasgo es mayor entre gemelos que entre mellizos (criados juntos, del mismo sexo), se sospecha componente gen\'etico:
\[
    C_G = 2\,[\mathrm{Corr}(\text{gemelos}) - \mathrm{Corr}(\text{mellizos})], \quad
    C_{AC} = \mathrm{Corr}(\text{gemelos}) - C_G, \quad
    C_{ANC} = 1 - C_G - C_{AC}.
\]
\textbf{Supuesto clave (\emph{equal environments assumption}):} los ambientes que generan similitudes son iguales para gemelos y mellizos. Problemas: si los gemelos comparten ambientes m\'as parecidos, se \textbf{sobreestima} $C_G$; si los ambientes intrauterinos son m\'as similares entre mellizos, se \textbf{subestima}. Adem\'as, problema de generalizaci\'on (m\'as prematuros que en la poblaci\'on general).

\subsubsection{Comparaci\'on con hijos adoptivos}

\begin{itemize}[nosep]
    \item \textbf{Padres/hijos criados juntos:} si la correlaci\'on biol\'ogicos $>$ adoptivos, hay componente gen\'etico. $C_G = 2\,[\mathrm{Corr}(\text{biol\'ogicos}) - \mathrm{Corr}(\text{adoptivos})]$; $C_{AC} = \mathrm{Corr}(\text{adoptivos})$.
    \item \textbf{Hijos dados en adopci\'on al nacer:} la similitud con padres \emph{biol\'ogicos} estima $C_G$; con padres \emph{adoptivos}, el componente ambiental. $C_G = \mathrm{Corr}(\text{biol\'ogicos dados en adopci\'on})$; $C_{AC} = \mathrm{Corr}(\text{adoptivos})$.
\end{itemize}
Problemas: \textbf{colocaci\'on selectiva} (agencias buscan familias parecidas a las biol\'ogicas $\to$ sobreestima $C_G$ y $C_{AC}$), ambiente intrauterino y generalizaci\'on.

\subsection{M\'etodos modernos: correlaci\'on entre alelos y rasgos}

Se miden variantes en los nucle\'otidos del ADN (\textbf{SNPs}, \emph{single nucleotide polymorphisms}) mediante estudios de asociaci\'on del genoma completo (\textbf{GWAS}). Ejemplo: \textit{Gene discovery and polygenic prediction from a GWAS of educational attainment} (Nature Genetics). Para inferir causalidad se controlan variables demogr\'aficas y se estudia el ``puente'' entre gen y comportamiento (e.g.\ genes de logro educativo se expresan en corteza prefrontal e hipocampo).

\textbf{Puntaje poligénico:} combina lectura de ADN con an\'alisis matem\'atico para estimar la probabilidad de que una persona tenga o desarrolle cierto rasgo. Incluye todos los SNPs que correlacionan con el rasgo (aunque individualmente no sean significativos, como en un test psicol\'ogico). El puntaje se basa en la presencia/ausencia de m\'ultiples SNPs.

\textbf{Heredabilidad faltante (\emph{missing heritability}):} diferencia entre la heredabilidad estimada por estudios familiares (gemelos) y la explicada por variantes gen\'eticas identificadas en GWAS.

\subsection{Correlaci\'on genotipo-ambiente}

La influencia gen\'etica puede ser \textbf{directa} o \textbf{indirecta} (mediada por la cultura): \emph{los genes influyen sobre los ambientes}. Tres tipos:
\begin{itemize}[nosep]
    \item \textbf{Evocativa:} los chicos son tratados seg\'un sus propensiones gen\'eticas (un chico con facilidad para la matem\'atica recibe m\'as regalos con n\'umeros).
    \item \textbf{Activa:} los chicos crean ambientes que correlacionan con sus propensiones (buscan juguetes con n\'umeros).
    \item \textbf{Pasiva:} reciben genes que correlacionan con los ambientes en que viven (padres con alelos que predisponen a la lectura $\to$ m\'as libros en casa).
\end{itemize}
\textbf{Consecuencias:} el componente gen\'etico puede ser indirecto y por tanto modificable; muchos ambientes que correlacionan con rasgos lo hacen porque los genes influyen en ambos (correlaci\'on pasiva). \emph{Ejemplo:} abandono paterno y depresi\'on de los hijos pueden no ser causales; controlando el puntaje poligénico de los hijos la relaci\'on desaparece (correlaci\'on genotipo-ambiente pasiva).

\subsection{Interacci\'on genotipo-ambiente}

Distinta de la correlaci\'on: la \textbf{influencia de los ambientes en los rasgos puede ser diferente para diferentes alelos} (o la influencia de los alelos, diferente seg\'un el ambiente). Ejemplos (dudosos): gen MAOA y maltrato infantil $\to$ comportamiento antisocial; gen transportador de serotonina y eventos estresantes $\to$ depresi\'on. Estudios recientes con puntaje poligénico a menudo \textbf{no} detectan interacci\'on (e.g.\ maltrato aumenta s\'intomas de TDAH independientemente del puntaje poligénico: rectas paralelas).

\end{document}
